\subsection{Nearest work: epistemic state graphs and recursive-reasoning termination}

The most direct prior art for coupling an explicit reasoning state to a stopping rule is Guha et al.'s 2026 framework for recursive reasoning systems \cite{guha2026termination}. Their epistemic state graph stores claims, evidential relations, conflicts, partial answers, open questions and confidence weights. They decompose iteration into an external-evidence expansion operator and an internal consolidation operator, then define an \emph{order-gap}: the distance between expand-then-consolidate and consolidate-then-expand. A windowed small order-gap is used as an endogenous stopping signal. Their main formal result characterizes when the linearized commutator near a consolidation fixed point is non-degenerate; the paper explicitly does not claim that a small order-gap proves global convergence or truth.

This result owns the general idea that recursive reasoning should expose both state and termination instead of relying on a token budget or iteration count. It also exposes a real weakness in a naive repeated-flatness rule: one observed operator ordering can look stable while a perturbation of the ordering still changes the epistemically meaningful state. The RAKL saturation contract therefore includes an explicit operator-order perturbation audit. It does not reuse the Euclidean order-gap as its certificate. Instead it runs the alternative orderings under the same frozen basis and computes their difference in the non-compensatory epistemic-growth coordinates. Different serialization or graph digests are harmless when the substantive difference vector is zero; any change to mechanisms, derivations, evidence roots, counterexamples, novelty, assumptions, fibers or discovery routes blocks saturation.

The distinction matters because the stopping questions are different. Guha et al. ask when a recursive reasoning trajectory has become locally insensitive to the order of expansion and consolidation. RAKL asks when a publishable scientific state has stopped changing under heterogeneous evidence/discovery/review perturbations while all authority, novelty, unresolved-fiber and freshness obligations remain closed. The former is an operator-dynamics criterion; the latter is a governance certificate over an evidence-bearing release state. RAKL now treats operator-order stability as one necessary coordinate of that larger certificate rather than as a competing termination theory.

A second adjacent 2026 proposal, Epistemic State Replication, separates an immutable evidence log from a stochastic belief lineage when distributed agent replicas need semantic rather than bitwise agreement \cite{heyu2026esr}. This reinforces the present separation between canonical evidential state and model-dependent computational state: semantic agreement between agents is not identical to evidence authority, but a deterministic evidence boundary can anchor divergent internal reasoning. The two systems address different execution problems, yet the shared separation is useful prior art and narrows any claim that evidence/belief decoupling is unique to RAKL.

Finally, a 2026 National Academies-associated framework for research trustworthiness argues for a systems view spanning accountability, evaluability, formulation, evaluation, bias control, error reduction and whether claims are warranted by evidence \cite{nosek2026trust}. This provides an external metascientific reason not to identify scientific trustworthiness with one confidence coordinate. The RAKL authority tuple is not asserted to be a canonical decomposition of trustworthiness; it is an executable local certificate structure whose multidimensional form is consistent with the broader finding that trustworthy science depends on several non-interchangeable properties.
